%% This is file `DEMO-RWTHLetter.tex' version 1.1 (2026-02-13),
%% it is part of
%% RWTH-CI -- Corporate Design for RWTH Aachen
%% ----------------------------------------------------------------------------
%%
%% Copyright (C) 2025–2026 by Marei Peischl <marei@peitex.de>
%%
%% ============================================================================
%% This work may be distributed and/or modified under the
%% conditions of the LaTeX Project Public License, either version 1.3c
%% of this license or (at your option) any later version.
%% The latest version of this license is in
%% http://www.latex-project.org/lppl.txt
%% and version 1.3c or later is part of all distributions of LaTeX
%% version 2008/05/04 or later.
%%
%% This work has the LPPL maintenance status `maintained'.
%%
%% The current maintainer of this work is
%%   Marei Peischl <rwth-ci@peitex.de>
%%
%% The development repository can be found at
%% https://gitlab.git.nrw/rwth-it-center/rwth-latex-templates/rwth-ci
%% Please use the issue tracker for feedback!
%%
%% If you need a compiled version of this document, have a look at
%% http://mirror.ctan.org/macros/latex/contrib/rwth-ci/doc
%% or at the documentation directory of this package (if installed)
%% <path to your LaTeX distribution>/doc/latex/rwth-ci
%% ============================================================================
%%
% !TeX program = lualatex
%%

\documentclass[
	english,%language as global option
	parskip=full,% select parskip mode
]{scrartcl}

\usepackage[
% textwidth=wide,% choose wide letter layout instead of narrow
% payment-data=true,% Enable the output of payment data, may be empty, and may be set in the lco file
% logofile=example-image-plain,% Select your logo file
% accept-missing-logos=true,% No error in case logo files are not available
]{rwth-letter}

%%%%%%%%%%%%%%%%%%%
% Language setup
%%%%%%%%%%%%%%%%%%%
\usepackage[english]{babel}
\usepackage[autostyle]{csquotes}% \enquote, to simplify use of quotation marks

\LoadLetterOption{DEMO-RWTHFromaddress}% Load address data from lco-file

\begin{document}

\begin{letter}{%
Mr. Lion\\
hbox 23/42\\
1337 Duck City\\
Island of \TeX
}

\setkomavar{subject}{\LaTeX{} letters using RWTH Aachen CI}
% \setkomavar{date}{2024-02-26}% In case \today should not be used
\setkomavar{yourmail}{yyyy-mm-dd}
\setkomavar{myref}{Demo-xx}

\opening{Hello,}
this is a template file to show the usage of the rwth-letter class.

The most important options can be found within the source file of this document.
Other options or layout adjustments may not comply CI regulations and therefore  should not be used.

Supported options:\\

\begin{description}
\item[raggedright=true/false] Set all text ragged right. Default setting is false and uses justified text.
\item[textwidth=narrow/wide] The design allows two different letter layouts. The default is narrow.
\end{description}

To reuse the sender's address the template provides an example lco-file.
See the \KOMAScript{} documentation on details about these files and how to use them.

\closing{Happy \TeX{}ing}

\vfil
\encl{Sources for this document:\\
\texttt{DEMO-RWTHLetter.tex} (Letter),\\
\texttt{DEMO-RWTHFromaddress.lco} (Address data)}

\end{letter}



\end{document}
%% End of file `DEMO-RWTHLetter.tex'.
